% Generated draft: requires mathematical and source review.
\documentclass[11pt]{article}
\usepackage[margin=25mm]{geometry}
\usepackage{fontspec}
\setmainfont{Latin Modern Roman}
\usepackage{amsmath,amssymb,mathtools,unicode-math}
\setmathfont{Latin Modern Math}
\usepackage{enumitem,needspace,xurl,hyperref}
\setlength{\parindent}{0pt}
\setlength{\parskip}{5pt}
\setlength{\emergencystretch}{3em}
\newcommand{\N}{\mathbb{N}}
\newcommand{\Z}{\mathbb{Z}}
\newcommand{\Q}{\mathbb{Q}}
\newcommand{\R}{\mathbb{R}}
\newcommand{\C}{\mathbb{C}}
\newcommand{\F}{\mathbb{F}}
\newcommand{\A}{\mathbb{A}}
\newcommand{\PP}{\mathbb{P}}
\newcommand{\E}{\mathbb{E}}
\newcommand{\eps}{\varepsilon}
\newcommand{\dd}{\,\mathrm{d}}
\begin{document}
\section*{UnsolvedMath ID: 20000527}
\textbf{Source problem code:} catalogue entry
\subsection{Definitions and mathematical statement}
This permanent record contains the entire thirteen-item workshop list, not a single Monge--Amp\`ere question. Fix a normalization of \(dd^c\); on a compact K\"ahler background \(\omega\), distinguish \((\omega+dd^cu)^n\) from the local expression \((dd^cu)^n\). The source does not consistently state backgrounds, solution classes, or convergence topologies. Its recovered problems are:
\begin{enumerate}
\item Monge--Amp\`ere mass on analytic sets (S. Dinew): seek a suitable definition for singular potentials, asking independence of bounded decreasing approximations with the same limit, and asking whether the weak limit for \(u_j=\max(u,-j)\) defines \(\operatorname{MA}(u)\). Seek \(\operatorname{MA}(u)=\mu\) when \(\mu\) is supported on an analytic set other than a point, possibly with prescribed Green-type singularities. Ask similar Dirichlet questions on \(\Omega\Subset\C^n\), real/toric analogues, and the following degeneration questions. For a Fano \(X\) with no KE metric, can normalized continuity-method potentials \(\psi_t=\varphi_t-\sup\varphi_t\) have a subsequential limit with defined MA measure supported on the associated multiplier-ideal locus (example: \(\PP^2\) blown up at one point)? For Fano \(X\) and smooth anticanonical \(D\), study a subsequential limit of \(\omega_\varepsilon\) satisfying \(\operatorname{Ric}(\omega_\varepsilon)=\varepsilon\omega_\varepsilon+(1-\varepsilon)[D]\), and whether a suitably defined limiting Ricci current is supported on \(D\) (H. Guenancia). The limiting parameter and singular-current conventions are unstated. The phrase ``compact K\"ahler case'' also gives \(\PP^n\setminus H\), which is not itself compact; the intended ambient/background setting needs clarification. Standard non-pluripolar MA products do not charge analytic pole sets, so they cannot be silently chosen as the requested charged-mass definition.
\item Find non-product smooth solutions of \((dd^cu)^n=0\) on a compact \(M\), for example \(\PP^2\), without assuming \(u\) psh (Y. Rubinstein). ``Non-product'' is undefined.
\item Find psh solutions of \((dd^cu)^n=0\) on \(\C^n\) (W. He). Constants and lower-variable examples make the literal question too weak without further growth, nontriviality, or exhaustion conditions.
\item On a strongly pseudoconvex \(\Omega\) with \(C^{3,1}\) boundary, seek an analytic proof independent of Krylov of \(u\in C^{1,1}(\Omega)\) for the intended Dirichlet solution \((dd^cu)^n=f\,dV\), \(f\geq0\), \(f^{1/n}\in C^{1,1}\), boundary value \(\varphi\in C^{3,1}\).
\item Ask the same with \(f^{1/(n-1)}\in C^{1,1}\), in dimension \(n\geq2\). The source says this is not covered by Krylov and records Pli\'s's affirmative special case for the ball with zero boundary value. These are historical qualifications, not a current complete status assessment.
\item Can Ross--Witt-Nystrom's HCMA examples without a foliation yield a counterexample to the maximal-rank question (M. Paun)? The exact maximal-rank statement is absent.
\item Solve \((\pi^*\omega+i\partial\overline\partial u)^{n+1}=0\) on \(X\times A\), \(\pi\) the projection and \(A\) an annulus, with boundary data \(\varphi_0,\varphi_1\) on the circles of radii \(1,e\), these data \(\omega\)-psh with unspecified regularity, \(\omega\) possibly degenerate, and \(\int_X\omega^n>0\) (E. Di Nezza).
\item Find a PDE proof of Ko\l odziej's \(L^\infty\) estimate and the optimal constant for a ball (Z. Blocki); the data and normalization are unstated.
\item Find a KE Fano \(X\) and \(u\in\Lambda_1\) with \(\int_Xu^3\omega_{\mathrm{KE}}^n\ne0\) (H. Macbeth). The eigenspace/Laplacian convention must be specified.
\item Does \((dd^cu)^n=1\) together with \(\Delta u\in L^{n(n-1)}\) imply \(u\in C^\infty\) (T. Collins)? The source omits the solution class and interior/boundary scope.
\item Seek a complex analogue of Pogorelov's estimate; no precise inequality is given.
\item For \(p:X\to D\), \(X\) K\"ahler and \(K_X\) nef, study solutions on \(X_t\) of \(\operatorname{Ric}(\omega_{t,\varepsilon})=-\omega_{t,\varepsilon}-\varepsilon\beta_t\) (M. Paun). The base, \(\beta_t\), degeneration, and desired estimates are unspecified.
\item Find an analytic proof of the ACC conjecture/theorem (T. Collins); the invariant and precise ACC statement are not identified.
\end{enumerate}

The report's explicit candidate concerns only item 9. On \(\PP^m\), \(m\geq2\), normalize the Fubini--Study KE form by \(\operatorname{Ric}(\omega)=\omega\), and use the first positive complex-Laplacian convention \(\Delta u=u\). It asserts that
\[
u([Z])=\frac{|Z_0|^2}{\sum_{j=0}^m|Z_j|^2}-\frac1{m+1}\in\Lambda_1,
\qquad
\frac{\int_{\PP^m}u^3\omega^m}{\int_{\PP^m}\omega^m}
=\frac{2m(m-1)}{(m+1)^3(m+2)(m+3)}>0.
\]
For \(m=2\) this equals \(1/135\). The report also claims a divisorial truncation formula relevant to item 1, a proper-exhaustion obstruction for item 3, and a conservative status map for all thirteen items; none of those formulas or hypotheses is supplied, so no complete additional result is reconstructed.
\subsection{Short English statement}
The record merges thirteen Monge--Amp\`ere workshop questions and must retain them all. The report's explicit projective-space moment formula addresses item 9 only; the many omitted definitions and hypotheses in the remaining items are recorded for review.
\subsection{Sources}
\begin{itemize}
\item \emph{Complex Monge--Amp\`ere Equation Workshop: Open problems}, edited by M. Dellatorre, September 8, 2016: \url{https://aimath.org/pastworkshops/mongeampereproblems.pdf}. The supplied verification says it is a two-page thirteen-item list; no PDF or report was accessed here.
\item Local canonical registry \texttt{aim-analysis-notes.json}, \texttt{notes[58]}, tag \texttt{problem}; permanent UnsolvedMath ID \(20000527\), source code \texttt{catalogue entry}. The extraction merges all items, makes item 5 ``15'' by joining a footer, loses the fractional exponents and \(\varepsilon\), and appends page 2. The clean statement is absent.
\item Source attributions: S. Dinew, H. Guenancia, Y. Rubinstein, W. He, M. Paun, E. Di Nezza, Z. Blocki, H. Macbeth, and T. Collins; source comparisons name Krylov, Pli\'s, Ross--Witt-Nystrom, Ko\l odziej, and Pogorelov. No full bibliographic entries or precise ACC reference are supplied. Original remarks and literature field are empty.
\item Supplied research attempt 1 proposes the projective-space family and mentions additional truncation/exhaustion results. It is machine generated, not peer reviewed, and requires independent verification before citation.
\end{itemize}
\end{document}
